\section{Discussion}
\label{sec:discussion}

\subsection{What a process team should take from this}

The most transferable result here is Section~\ref{sec:interface}: analog level placement in
this device is set by fixed interface charge and essentially not by ferroelectric
thickness. A \SI{0.3}{\nano\metre} thickness tolerance moves the levels by \num{0.07}
decades; \SI{10}{\percent} of interface charge moves them by \num{1.31}. Combined with the
placement specification in Section~\ref{sec:variability} (worst-level misplacement inside
about \num{0.6} decades if only a gain trim is affordable), this converts a device-level
observation into a control target.

This ordering was not expected. Ferroelectric thickness sets the coercive voltage, so
thickness control looks like the obvious lever. A study that measured only thickness
sensitivity would have concluded the device was well-behaved and missed the dominant term.

\subsection{Three figures of merit that do not work}

This paper reports three separate measurements of level distinguishability and none of them
predicts network accuracy. That is a negative result about a widely used class of metric,
and it has a clear mechanism: separability asks whether a level can be resolved from its
neighbour \emph{at the readout}, while a network needs the realized weight to sit close to
its target. Those are different questions, and only the second one moves a classifier.

The practical consequence is a caution about design rules that emit integers. The
three-sigma separability criterion returns a usable level count, and following it here
would have cost \num{0.036} accuracy. Worse, the count itself is unstable against the data
it is computed from: eleven repeats give eight levels and a \num{0.036} penalty, while four
repeats of the same run give six levels and a penalty ten times larger. A rule whose output
is an integer derived from a threshold test on a quantity known to \SI{22}{\percent} is not
a safe basis for a hardware decision.

\subsection{Why the level-count question has two answers}

The device supports eight levels open-loop and about sixteen with write-verify, and this
paper reports both with their protocols attached. The distinction is not a hedge; it is
the result. The train's own stops bunch at the top of a saturating curve, below the
cycle-to-cycle floor, so counting pulses undercounts what the range can hold; placing
targets across the same range with verification uses it fully.

Presenting only the larger number would be an overclaim, and presenting only the smaller
one would understate the device by a factor of two. The reason the paper is comfortable
deploying on fifteen levels is not that fifteen is defensible in isolation but that
Section~\ref{sec:variability} shows write-verify is required regardless: without per-device
calibration, accuracy on the corner devices is \num{0.243}. The system pays for verify
circuitry for a reason that has nothing to do with level count, and once it is paid the
higher count is free.

\subsection{The system, not the device}

Section~\ref{sec:energy} is the part of this work most likely to be quoted out of context,
so it is stated plainly here as well. The computing core costs
\SI{535.8}{\pico\joule} per heartbeat. The analog-to-digital conversion a real chip would
need, under a standard \SI{1}{\pico\joule} per conversion assumption, comes to
\SI{183}{\nano\joule}, \num{342} times more. Device-level energy figures for
compute-in-memory arrays are frequently reported without this term, and on this workload
that omission would change the headline number by more than two orders of magnitude.

This does not make the device uninteresting. Programming the whole array costs
\num{0.13} of one inference and is paid once, which is a real and measured argument for
non-volatile analog weights. It does mean that further work on this device's write energy
has little system-level value until the conversion overhead is addressed, and that is more
useful for a reader to know than a smaller headline number would have been.

\subsection{Limitations}
\label{sec:limits}

\paragraph*{The model is two-dimensional} The device is a double-gate cross-section mapped
to a nanosheet through the gate perimeter. The mapping was verified to be exact arithmetic
to double-precision epsilon, and that is the full extent of what was verified. No
three-dimensional gate-all-around simulation was run, so corner electrostatics and
edge fields along the sheet width are not represented. This is expected to matter most for
absolute current and least for the ratios and voltages this paper reports, but that
expectation is an argument rather than a measurement.

\paragraph*{The absolute current level is not calibrated} The calibration fits a free width
scalar which absorbs a factor of 141 between the simulated and measured saturation
currents. The window, threshold voltages, subthreshold slope, on/off ratio and turn-on
shape are calibrated; the current density is not, and no conclusion here depends on it.

\paragraph*{Endurance is not measured and cannot be} The Preisach ferroelectric model has
no fatigue, wake-up or imprint term. Cycling it produces a flat curve because the equations
say so, not because the device is durable. The three-cycle data in Section~\ref{sec:device}
is write repeatability. Endurance on this stack is a measurement, and it has to be made on
silicon.

\paragraph*{Retention is bounded, not extrapolated} The hold spans
\SI{10}{\milli\second}. Ten years is eleven decades beyond that, and the plateau is flat to
the solver's resolution, so no ten-year figure is given. No activation energy is extracted
either: both temperatures reach flat plateaus and the model contains no thermally activated
retention term.

\paragraph*{The cycle-to-cycle statistic is from eleven repeats} The run was stopped by a
job time limit at eleven of twenty repeats, so the standard deviation carries about
\SI{22}{\percent} relative uncertainty. The claim made from it is correspondingly bounded:
at and around the measured noise, cycle-to-cycle variation costs essentially nothing. A
completed run would tighten the number but could not change that conclusion unless the true
value were roughly five times the measured one.

\paragraph*{Single-shot firing, not a leaky cycle} With the depolarization time constant
used for the application parameter set, the device holds its state after firing rather than
relaxing back. What Section~\ref{sec:device} demonstrates is a single integrate-and-fire
event. A repeating leaky cycle would need a different parameter set and is not claimed.

\paragraph*{One task, one architecture} The network results are for a single spiking
recurrent classifier on single-lead electrocardiogram beats, deployed by post-training
quantization from one fixed checkpoint. The placement-over-count conclusion follows from
the differential encoding on a four-decade ladder, which is general to that encoding, but
it has been demonstrated on one workload.

\subsection{On method}

Nine errors were found and corrected during this work, and several of them would have
produced a publishable-looking wrong result rather than an obvious failure. Three separate
measurements (retention, read disturb and cycling) initially showed dramatic
degradation that was entirely the post-write depolarization transient, misread as data
loss. Each looked like responsible bad news, which is exactly the kind of result that
survives review.

The rules in Section~\ref{sec:rules} are stated in the paper rather than kept in a
laboratory notebook because they are the difference between those numbers and these. Two
rules matter most: form a ratio only from measurements taken in the same run, and treat a
simulation that completes without an error as an unverified claim rather than a result. In
this work two runs completed cleanly with effectively no field anywhere, and 224 files were
produced by sweeps whose write pulses were too short to program the device at any
amplitude. All of them looked fine.
